% Table 3.1: Summary of Balanced Midwestern County Panel (1951-2025)
\begin{table}[htbp]
\centering
\small
\setlength{\tabcolsep}{8pt}
\caption{Descriptive summary of the 135 balanced study counties across the six Midwestern states (1951--2025).}
\label{tab:panel_selection}
\begin{tabular*}{\textwidth}{@{\extracolsep{\fill}}l c c c c}
\toprule
\textbf{State} & \textbf{Counties} & \textbf{Observations} & \textbf{Mean Yield} & \textbf{Std. Dev.} \\
 & ($N$) & ($N \times 75$) & (bu/ac) & (bu/ac) \\
\midrule
Illinois  & 28  & 2,100  & 38.35 & 13.20 \\
Indiana   & 20  & 1,500  & 39.52 & 11.82 \\
Iowa      & 35  & 2,625  & 39.73 & 12.01 \\
Minnesota & 19  & 1,425  & 35.10 & 12.70 \\
Missouri  & 6   & 450    & 29.04 & 10.54 \\
Ohio      & 27  & 2,025  & 36.88 & 11.40 \\
\midrule
\textbf{Panel Total} & \textbf{135} & \textbf{10,125} & \textbf{37.72} & \textbf{12.40} \\
\bottomrule
\end{tabular*}
\vspace{1ex}
\raggedright
\footnotesize{\textit{Notes:} Counties ($N$) indicates the number of counties selected across the six Corn Belt states based on unbroken 75-year continuity (1951--2025) under a strict zero-imputation protocol. Observations denotes total finite county-year yield records. Mean Yield and Std. Dev. are computed over the full 75-year panel.}
\end{table}
